\documentclass[11pt]{article}
\usepackage{amsmath,amssymb}
\usepackage{xurl}
\usepackage[colorlinks=true,linkcolor=blue,citecolor=blue,urlcolor=blue]{hyperref}
% Shared commands for notes in docs/paper-gaps/.

\DeclareUnicodeCharacter{2080}{\ifmmode{}_{0}\else\textsubscript{0}\fi}
\DeclareUnicodeCharacter{2081}{\ifmmode{}_{1}\else\textsubscript{1}\fi}
\DeclareUnicodeCharacter{2082}{\ifmmode{}_{2}\else\textsubscript{2}\fi}
\DeclareUnicodeCharacter{2099}{\ifmmode{}_{n}\else\textsubscript{n}\fi}

\newcommand{\Lean}{\textsc{Lean}}
\ExplSyntaxOn
\NewDocumentCommand{\leanid}{m}{
  \ifmmode
    \textsf{\detokenize{#1}}
  \else
    \group_begin:
    \urlstyle{sf}
    \tl_set:Nn \l_tmpa_tl {#1}
    \tl_replace_all:Nnn \l_tmpa_tl {\_} {_}
    \exp_args:NV \path \l_tmpa_tl
    \group_end:
  \fi
}
\ExplSyntaxOff
\providecommand{\tr}{\operatorname{tr}}
\newcommand{\tnleanIssueBase}{https://github.com/LionSR/TNLean/issues/}
\newcommand{\tnleanPrBase}{https://github.com/LionSR/TNLean/pull/}
\newcommand{\tnleanDocBase}{https://sirui-lu.com/TNLean/docs/}
\newcommand{\ghissue}[1]{\href{\tnleanIssueBase#1}{GitHub issue \##1}}
\newcommand{\ghpr}[1]{\href{\tnleanPrBase#1}{PR \##1}}
\newcommand{\doclink}[1]{\href{\tnleanDocBase#1.html}{\path{#1}}}

% Dirac notation and the identity operator, as in blueprint/src/macros/common.tex.
% \providecommand keeps note-local definitions authoritative.
\usepackage{dsfont}
\providecommand{\ket}[1]{|#1\rangle}
\providecommand{\bra}[1]{\langle #1|}
\providecommand{\braket}[2]{\langle #1 | #2 \rangle}
\providecommand{\ketbra}[2]{|#1\rangle\!\langle #2|}
\providecommand{\Id}{\mathds{1}}
\providecommand{\one}{\mathds{1}}
\providecommand{\C}{\mathbb{C}}
\providecommand{\MN}[1]{M_{#1}(\C)}
\providecommand{\MD}{M_D(\C)}

% Verdict marker: \gapnote{<kind>}{<status>}. Kind is the policy
% classification (clarification, local-correction, scope-restriction,
% unfaithful, false-source, open-gap); status is open, wip (elimination
% underway), resolved, or historical. Machine-read by the site index;
% prints a verdict line.
\newcommand{\gapnote}[2]{\par\noindent\textbf{Verdict:} #1 (#2).\par}

\title{Scalar Phases in the Physical String-Order Criterion}
\author{}
\date{2026-10-05}
\begin{document}
\maketitle
% Full production counterexample target checked at 4e22a416, PR CI
% 37385801777, job 112019128398, step 5. Full-root/regression gates remain separate.
\gapnote{false-source}{resolved}

\section{The two literal assertions}

P\'erez-Garc\'ia, Wolf, Sanz, Verstraete, and Cirac define string order by
\begin{align}
 S_N(x,y,u)&=\langle x\otimes u^{\otimes N}\otimes y\rangle,
 &\lim_{N\to\infty}|S_N(x,y,u)|&>0,
 \label{eq:scalar_phase_source_order}
\end{align}
for some local unitary $u\ne\Id$ and local endpoint operators $x,y$, which
may be taken Hermitian. The exclusion is the literal identity matrix,
not every scalar unitary. This is display SOP, lines 114--121 of the local
source for arXiv:0802.0447 \cite{PerezGarcia2008StringOrder}.

For a pure finitely correlated state, Theorem 1 asserts that this property
is equivalent to the existence of a unitary $\widetilde u\ne\Id$, a virtual
matrix $V$, and physical letters $n,m$ such that
\begin{align}
 \mathcal E_{\widetilde u}(V)&=V,
 &\operatorname{tr}(V\Lambda A_n A_m^\dagger)&\ne0.
 \label{eq:scalar_phase_source_fixed}
\end{align}
The source passage preceding this theorem replaces the original twist $u$
by $\widetilde u=\mu^{-1}u$, where $\mu$ is its peripheral eigenvalue.
Although this preserves unitarity, it need not preserve $u\ne\Id$:
if $u=\mu\Id$, the new twist is the identity. The relevant source passages
are lines 257--264 and 270--275.\footnote{Local source:
\path{Papers/0802.0447/StringOrder-v10.tex}.}

\section{The literal string-order predicate is automatic}

For every normalized state with positive physical dimension, choose
$x=y=\Id$ and $u=-\Id$. Then
\begin{align}
 S_N(\Id,\Id,-\Id)&=(-1)^N,
 &|S_N(\Id,\Id,-\Id)|&=1.
 \label{eq:scalar_phase_trivial_order}
\end{align}
Both endpoints are Hermitian. For a canonical tensor this is the exact
transfer calculation
$\mathcal E_{\mu\Id}=\mu\mathcal E$,
$\mathcal E(\Id)=\Id$, and $\operatorname{tr}\Lambda=1$.
Thus the literal definition supplies no nontriviality test on the state.

This observation alone does not disprove the existential conclusion in
(\ref{eq:scalar_phase_source_fixed}): a state could possess some other
nonidentity fixed twist. The following example rules out all such twists.

\section{A nondegenerate canonical example}

Write $\sigma_x,\sigma_y,\sigma_z$ for the Pauli matrices. For
$k=x,y,z$, set
\begin{align}
 U_0&=\Id,
 &U_k&=\frac{\Id-i\sigma_k}{\sqrt2},
 \label{eq:scalar_phase_tensor}\\
 A_j&=\sqrt{p_j}\,U_j,
 &\Lambda&=\frac{\Id}{2},
 \notag\\
 (p_0,p_x,p_y,p_z)&=\left(\frac12,\frac14,\frac16,\frac1{12}\right).
 \notag
\end{align}
Each $U_j$ is unitary. Hence the transfer map is unital and trace
preserving, and $\Lambda$ is faithful, stationary, and has trace one.
The four Kraus matrices are linearly independent: after removing their
nonzero scalar factors, they are
$\Id,\Id-i\sigma_x,\Id-i\sigma_y,\Id-i\sigma_z$, an invertible triangular
change of the Pauli basis. Therefore they span $M_2(\C)$. For each nonzero
vector $v$, the vectors $A_jv$ span $\C^2$, so
$\mathcal E(|v\rangle\langle v|)$ is positive definite. Full one-letter
span consequently implies irreducibility and primitivity. In particular, every nonzero
matrix $X$ satisfying
$\mathcal E(X)=\nu X$ with $|\nu|=1$ has $\nu=1$ and is scalar.
The example therefore meets the source's canonical pure-FCS hypotheses;
its bond space and physical alphabet are both nondegenerate.

On the traceless Pauli basis the transfer action is
\begin{align}
 T&=\begin{pmatrix}
 3/4&-1/12&1/6\\
 1/12&2/3&-1/4\\
 -1/6&1/4&7/12
 \end{pmatrix}.
 \label{eq:scalar_phase_bloch}
\end{align}
Suppose a unitary $u$ and nonzero $Q$ satisfy $\mathcal E_u(Q)=Q$.
The peripheral rigidity statement of the source's Lemma 1 supplies a
unitary $W$ and $|\mu|=1$ with
$\sum_j u_{ij}A_j=\mu W A_iW^\dagger$.
Peripheral uniqueness forces $\mu=1$. Consequently the transfer map
commutes with conjugation by $W$.

Let $R$ be the complex $3\times3$ matrix of this conjugation on the Pauli
basis. Trace preservation of the bilinear form
$(X,Y)\mapsto\tfrac12\operatorname{tr}(XY)$ gives $R^{\mathsf T}R=\Id$.
Covariance gives $RT=TR$. Transpose this relation and use orthogonality
to obtain $RT^{\mathsf T}=T^{\mathsf T}R$. Thus $R$ commutes with
\begin{align}
 T+T^{\mathsf T}&=\operatorname{diag}(3/2,4/3,7/6).
 \notag
\end{align}
Its diagonal entries are distinct, so $R$ is diagonal. All off-diagonal
entries of $T$ are nonzero, so $RT=TR$ makes all three diagonal entries of
$R$ equal. Write $R=c\Id$. Orthogonality gives $c^2=1$, whereas preservation
of $\sigma_x\sigma_y=i\sigma_z$ gives $c^2=c$. Therefore $c=1$.

Conjugation by $W$ fixes the identity and all Pauli matrices, hence every
matrix. The intertwining equations now read
$\sum_j u_{ij}A_j=A_i$. Linear independence of the physical letters yields
$u=\Id$. This contradicts the required nonidentity twist in
(\ref{eq:scalar_phase_source_fixed}); its nonzero trace coefficient already
implies $Q\ne0$. Yet (\ref{eq:scalar_phase_trivial_order}) proves string
order for this same tensor. The literal existential Theorem 1 is therefore
false as printed.\footnote{The explicit declarations are
\leanid{MPSTensor.stringPhaseTensor_canonical},
\leanid{MPSTensor.stringPhaseTensor_pure}, and
\leanid{MPSTensor.stringPhaseTensor_literal_criterion_counterexample},
in \path{TNLean/MPS/Examples/StringOrderScalarPhase.lean}.}

\section{A phase-retaining correction}

Keep the original physical twist $u$. Under the canonical pure-FCS
hypotheses, physical endpoints with positive limiting magnitude exist
if and only if there are a virtual matrix $V$, a phase $|\mu|=1$, and
physical letters $n,m$ such that
\begin{align}
 \mathcal E_u(V)&=\mu V,
 &\operatorname{tr}(V\Lambda A_nA_m^\dagger)&\ne0.
 \label{eq:scalar_phase_corrected}
\end{align}
To prove necessity, positive limiting magnitude excludes twisted spectral
radius less than one. The unitary-twist bound and peripheral rigidity then
supply a virtual unitary intertwiner. The physical endpoint coefficient
criterion supplies the two letters. Conversely the nonzero coefficient
makes $V$ nonzero. Peripheral rigidity and uniqueness normalize its
one-dimensional eigendirection to a unitary intertwiner, and the same
endpoint criterion constructs physical endpoints.

\section{The explicit projective convention}

For physical string order, require the unitary twist to satisfy
\begin{align}
 u&\notin\{c\Id:c\in\C\}.
 \label{eq:scalar_phase_nonscalar}
\end{align}
This is an explicit correction to the printed definition, rather than an
unstated assumption on the source theorem. If $|\mu|=1$, then $u$ satisfies
(\ref{eq:scalar_phase_nonscalar}) exactly when $\mu^{-1}u$ does: an equality
$\mu^{-1}u=c\Id$ would give $u=(\mu c)\Id$.
Consequently the phase-retaining criterion implies the corrected global
fixed-point criterion. Physical string order in this projective sense is
equivalent to the existence of a nonscalar unitary $\widetilde u$, a virtual
matrix $V$, and letters $n,m$ satisfying
\begin{align}
 \mathcal E_{\widetilde u}(V)&=V,
 &\operatorname{tr}(V\Lambda A_nA_m^\dagger)&\ne0.
 \label{eq:scalar_phase_projective_fixed}
\end{align}
The reverse direction follows from (\ref{eq:scalar_phase_corrected}) with
$\mu=1$.\footnote{The definition
\leanid{MPSTensor.HasPhysicalStringOrder} now states the nonscalar condition.
The global criterion is
\leanid{MPSTensor.hasPhysicalStringOrder_iff_exists_fixed_letter}.
The fixed-twist criterion is
\leanid{MPSTensor.pureCanonical_exists_physicalStringOrderWith_iff_peripheral_letter},
both in \path{TNLean/MPS/Symmetry/PhysicalStringPhase.lean}.}

The four-letter example has no physical string order in this corrected
sense, while its literal nonidentity-twist condition remains true.
Both the AKLT and two-site cluster physical witnesses are nonscalar, so their
physical endpoint conclusions are unchanged. Only this corrected physical
predicate is retained. The literal condition appears in the explicit
counterexample and scalar-twist calculation; it is not a second general
predicate. The fixed-twist physical condition remains unchanged, as does
virtual-boundary nondecay for fixed twists, including the identity.

The finite-support comparison and the $D^2$ Hermitian-product endpoint bound
are established separately in
\path{pgwsvc08_string_order_virtual_boundary.tex}, using
\leanid{MPSTensor.pureCanonical_exists_isHermitian_physicalStringProductOrder}.
The physical reduced-state symmetry theorem is a separate source assertion.

\bibliographystyle{alpha}
\bibliography{references}
\end{document}
